\section{Análisis de Seguridad y Comparativa}

Como parte de este estudio, nos toca hablar sobre qué tan seguro es realmente el estándar ruso GOST R 34.10 y cómo se le para enfrente al ya conocido ECDSA occidental que vemos por todos lados.

\subsection{Fortaleza matemática y análisis de seguridad}

Lo primero que hay que entender es que toda la seguridad del algoritmo GOST R 34.10--2012 recae sobre la dificultad del \textbf{Problema del Logaritmo Discreto en Curvas Elípticas (ECDLP)}. En español: aunque un atacante sepa qué curva estás usando (los parámetros $p, a, b, q$), conozca el punto generador $P$ y tu clave pública $Q$, averiguar tu clave privada $d$ requeriría resolver la ecuación $Q = d \times P$. Con los tamaños de clave actuales (256 o 512 bits) sobre cuerpos finitos, las computadoras más potentes de hoy en día tardarían miles o millones de años en resolverlo usando algoritmos conocidos (como Pollard’s rho o Baby-step giant-step).

Un punto fuerte a destacar aquí, especialmente en la forma en la que se calcula la firma, es la protección inherente contra los famosos \textit{timing attacks} (ataques de tiempo). Cuando generamos nuestro par de firma $(r, s)$ en el estándar GOST, la ecuación para $s$ es:
$$ s = (r \cdot d + k \cdot e) \pmod q $$
donde $d$ es la clave privada, $k$ es el número aleatorio efímero, y $e$ es el hash del mensaje convertido. 
Si lo analizamos bien, el algoritmo emisor del GOST \textbf{no requiere calcular el inverso modular del número aleatorio $k$} en la firma. Calcular un inverso modular es una operación matemáticamente "pesada" y su tiempo de ejecución puede variar mucho, lo que a menudo "filtra" información en implementaciones no cuidadosas. Que GOST lo omita del lado del emisor es un plus enorme para la seguridad en dispositivos limitados (como smart cards).

Además, el algoritmo está casado con la función hash rusa \textbf{Streebog} (GOST R 34.11-2012), la cual está diseñada para ser ultra resistente a ataques de colisión y de preimagen, superando en algunos análisis estructurales a las familias SHA convencionales gracias a su esquema basado en la red de sustitución y permutación (tipo AES/Kuznyechik).

\subsection{Comparativa técnica directa: GOST vs. ECDSA}

Aunque ambos estándares son "primos" porque operan bajo criptografía de curvas elípticas, tienen diferencias técnicas e ideológicas muy marcadas que son muy interesantes de ver en conjunto:

\begin{itemize}
    \item \textbf{Cálculo de la firma (Eficiencia):} \\
    En el estándar americano (ECDSA), generar $s$ requiere resolver: 
    $$s = k^{-1}(e + r \cdot d) \pmod n$$
    Como mencionamos, esto obliga al dispositivo del emisor a procesar un inverso modular ($k^{-1}$). En GOST, el emisor se libra de esta carga, trasladando la necesidad aritmética "fuerte" al receptor durante la verificación (quien normalmente tiene más poder de cómputo, por ejemplo, un servidor).

    \item \textbf{Las familias de familias de curvas:} \\
    ECDSA se suele utilizar con las curvas estándar del NIST (como la P-256). Durante años ha existido el debate y la sospecha en la comunidad criptográfica de que estas curvas podrían tener debilidades o "backdoors" introducidas a propósito por agencias de Estados Unidos. GOST utiliza sus propios conjuntos de los parámetros (como \texttt{id-tc26-gost-3410-2012-256}), lo cual le da a los gobiernos e instituciones que desconfían del NIST una alternativa matemáticamente robusta y libre de supuestas injerencias occidentales.

    \item \textbf{El manejo del Hash (Endianness):} \\
    Algo súper peculiar que notamos al implementar GOST (y especialmente la conexión con Streebog) es cómo manejan el formato numérico del hash. Mientras que ECDSA / SHA procesan sus bloques y salidas en formato \textit{Big-Endian}, Streebog y GOST operan estrictamente en \textit{Little-Endian}. Desde el punto de vista del desarrollo, esto requiere un módulo de transformación adicional para acomodar los bytes de Streebog antes de someterlos a las ecuaciones matemáticas en la curva, algo que ECDSA nos ahorra.
\end{itemize}

En conclusión, GOST R 34.10 no tiene absolutamente nada de desventaja frente a ECDSA. De hecho, su diseño asimétrico en cuanto a la carga computacional (hacerle la vida fácil al que firma y más compleja al que verifica) lo hace brillante en arquitecturas donde el cliente tiene poca memoria y procesamiento.